According to the Global Energy Review 2025 report published in March 2025 by the 
IEA \cite{international2025global}, global $\text{CO}_2$ emissions increased by 0.8\% in 2024, 
reaching an all-time high of 37.8 Gt $\text{CO}_2$ per year. In addition, the report highlights 
that energy demand grew by 2.2\% in 2024. This increase in $\text{CO}_2$ emissions is inconsistent 
with the goals of the Paris Agreement, which aims to reduce $\text{CO}_2$ emissions by 43\% by 2030
 and to reach net-zero emissions by 2050. However, the lower increase in $\text{CO}_2$ emissions
  compared to energy demand growth is an encouraging trend, as it suggests a relative decoupling 
  of emissions from energy consumption. \\

Kosmadakis et al. report that the waste heat potential in the EU-27
 and the UK reached 221.32 TWh/year in 2021 \cite{kosmadakis2024industrial}. In comparison, the demand for heat below 140°C 
 in Germany was 612 PJ, or 170 TWh, in 2012 \cite{arpagaus2018high}. These numbers justify the
 interest of companies and researchers in designing and studying machines capable of recovering
 this waste energy. \\

In this work, we will focus on the use of high-temperature heat pumps (HTHPs) for 
recovering waste heat from industrial processes. Arpagaus et al.  
provide a list of manufacturers and universities that have already been working on the design 
of several machines for at least one decade \cite{arpagaus2018high}. The existing machines cover a large spectrum of 
applications, ranging from small units of a few kilowatts to large ones of more than one megawatt. 
However, new design ideas can still be studied in order to improve current recovery systems.\\

This thesis focuses on the design, construction, and performance evaluation of a 
modular machine capable of operating under various thermodynamic cycles. 
Among these, the dual-evaporator transcritical cycle, whose schematic is shown in Figure 
\ref{fig:schematic}, represents the most elaborate and complex configuration 
implemented. A review of the existing literature did not reveal any previous applications 
of such a dual-evaporator transcritical cycle for HTHP systems. Furthermore, the machine developed 
within the framework of this master's thesis is intended to serve as a platform for 
future academic and research activities.